Whatever the case, the history of the differential and integral calculus, starting with the end of the XVIIth century, is divided into two parts. One is concerned with the applications of this calculus, always richer, more numerous and more varied. To the differential geometry of plane curves, to differential equations, to power series, to the calculus of variations, which have already been referred to above, have been joined the differential geometry of ill-shaped curves, then of surfaces, multiple integrals, partial differential equations, trigonometric series, the study of numerous special functions, and many other types of problems. We only should be concerned here with work which has contributed to clearing up, deepening and consolidating the very principles of the infinitesimal calculus, in what concerns functions of a real variable.

From this point of view, the great treatises of the middle of the XVIIIth century offer few novelties. Maclaurin in England \(^{17}\) {[}214{]}, Euler on the continent ({[}108 a{]}, (1), vv. X to XIII), remained faithful to the traditions of which each one was the inheritor. It is true that the first tries hard to clarify somewhat the Newtonian concepts \(^{18}\) whereas the second, pushing Leibnizian formalism to an extreme, is content, like Leibniz and Taylor, to base the differential calculus on a very obscure passage to the limit starting with the calculus of differences, a calculus of which as it happens he gives a very polished exposition. But above all Euler finishes the work of Leibniz by introducing and ensuring the adoption of the notation still in use today for e, i and the trigonometric functions, and spreading the notation π. On the other hand, and if most of the time he does not make any distinction between functions and analytic expressions, he insists, in connection with trigonometric series and the problem of vibrating strings, on the need not to limit oneself to the functions thus defined (and that he qualifies as ``continuous'') but to consider also, if the need arises, arbitrary, or ``discontinuous'', functions given experimentally by one or more arcs of a curve ({[}108 a{]}, (1), v. XXIII, pp.~74-91). Finally, even though this is somewhat outside our framework, it is not possible not to mention here his extension of the exponential function to the complex domain, from which he obtains the famous formulae linking the exponential to the trigonometric functions, as well as the definition of the logarithm of complex numbers; from that is elucidated definitively the famous analogy between the logarithm and the inverse circular functions, or, in the language of the XVIIth century, between the quadratures of the circle and of the hyperbola, observed already by Grégoire de St.~Vincent, made precise by Huygens and above all by J. Gregory, and which, with Leibniz and Bernoulli, appeared in the formal integration of

\[
\frac {1}{1 + x ^ {2}} = \frac {i}{2 (x + i)} - \frac {i}{2 (x - i)}.
\]

Meanwhile d'Alembert, the enemy of all mysticism in mathematics as elsewhere, had, in remarkable articles ({[}75 a{]}, articles DIFFÉRENTIEL and LIMITE, and {[}75 b{]}), defined with the greatest clarity the notions of limit and of derivative, and maintained strongly that in the end that is all the ``metaphysics'' of the infinitesimal calculus. But this wise advice did not have any immediate effect. The monumental work of Lagrange ({[}191{]}, vv. IX-X) represents an attempt to base analysis on one of the most arguable of the Newtonian concepts, that which confuses the notions of arbitrary functions and that of functions which can be expanded in a power series, and to draw from that (by consideration of the coefficient of the term of the first order in the series) the notion of differentiation. Of course, a mathematician of the quality of Lagrange could not fail to obtain in this case important and useful results, as for example (and in a way which was in fact independent of the point of departure that we have just indicated) the general proof of the formula of Taylor with the expression for the remainder in the form of an integral, and its evaluation by the theorem of the mean; in any case the work of Lagrange is at the origin of the method of Weierstrass in the theory of functions of a complex variable, as well as the modern algebraic theory of formal series. But, from the point of view of its immediate objective, it represents a retreat rather than progress.

With the teaching works of Cauchy, on the contrary ({[}56 a{]}, (2), v. IV) we find ourselves at last on solid ground. He defines a function essentially as we do it today, even though in a language which is still rather vague. The notion of limit, fixed once and for all, is taken as the point of departure; those of continuous function (in the modern meaning) and of derivative are immediately deduced, as well as their principal elementary properties; and the existence of the derivative, instead of being an act of faith, becomes a matter of study by the ordinary methods of analysis. Cauchy truth to tell, is hardly interested in this; and on the other hand, if Bolzano, having reached on his part the same principles, constructed an example of a continuous function not having a derivative at any point {[}27 a{]}, this example was not published, and the question was only publicly settled by Weierstrass, in a work in 1872 ({[}329 a{]}, v. II, pp.~71-74).

As far as integration is concerned, the work of Cauchy represents a return to the sane traditions of antiquity and the first part of the XVIIth century, but relying on technical means which are still insufficient. The definite integral, passed over into second place for too long, becomes again the primary notion, for which Cauchy ensures the definitive adoption of the notation \(\int_{a}^{b} f(x) dx\) proposed by Fourier (instead of the awkward \(\int f(x) dx\left[\begin{array}{ccc}x & = & a \\ x & = & b\end{array}\right]\) sometimes used by Euler); and, in order to define it, Cauchy returns to the method of exhaustion, or as we would say, to ``Riemann sums'' (which would be better called sums of Archimedes, or sums of Eudoxus). It is true that the XVIIth century did not deem it necessary to submit the notion of area to a critical examination, an idea which had appeared to them to be at least as clear as that of incommensurable real number; but the convergence of ``Riemann'' sums to the area under the curve, so long as it is the case of a monotonic curve or a piecewise monotonic curve, was a familiar idea to all the authors concerned with rigour in the XVIIth century, such as Fermat, Pascal, Barrow; and J. Gregory, particularly well prepared by his thoughts on the passage to the limit and his familiarity with an already very abstract form of the principle of ``nested intervals'', had even drawn up, it would appear, a careful proof, which remained unpublished ({[}136 d{]}, pp.~445-446), which could have been of use to Cauchy almost without any change if he had known it. \(^{19}\) Unfortunately for him, Cauchy claimed to prove the existence of the integral, that is to say the convergence of the ``Riemann sums'', for an arbitrary continuous function; and his proof, which would have become correct if it were based on the theorem on uniform continuity of functions continuous in a closed interval, is stripped of all convincing value for lack of this notion. Dirichlet does not seem to have seen either the difficulty at the time when he drew up his famous memoirs on trigonometric series, since he quotes there the theorem in question as ``easy to prove'' ({[}92{]}, v. I, p.~136); it is true that he only applies it definitively to bounded piecewise monotonic functions; Riemann, more circumspect, only mentions these latter when it is a matter of using his necessary and sufficient condition for the convergence of ``Riemann sums'' ({[}259 a{]}, pp.~227-271). Once the theorem on uniform continuity had been established by Heine (cf.~p.~145), the matter did not offer any more difficulties of course; and it is easily dealt with by Darboux in 1875 in his memoir on the integration of discontinuous functions {[}77{]}, a memoir in which he meets on many points with the important researches of P. du Bois-Reymond, which appeared at about the same time. In the same way for the first time, but this time definitively, the linearity of the integral of continuous functions is proved. On the other hand, the notion of uniform convergence of a sequence or series, introduced among others by Seidel in 1848, and put into prominence in particular by Weierstrass (cf.~p.~205) had allowed the construction of a solid basis, under conditions that were a bit too restrictive it is true, for the integration of series term by term, and for differentiation under the \(\int\) sign, while awaiting modern theories of which we do not have to speak here, and which would clarify these questions in a manner which is provisionally definitive.

We have thus attained the final stage of the classical infinitesimal calculus, that which is represented by the great Treatises on Analysis of the end of the XIXth century; from the point of view that concerns us, that of Jordan {[}174c{]} occupies amongst them a pre-eminent place, for aesthetic reasons on the one hand, but also because, if it constitutes an admirable setting out of the results of classical analysis, it announces in many ways modern analysis and prepares the way for it.

\chapter{ASYMPTOTIC EXPANSIONS.}

The distinction between the ``infinitely small'' (or ``infinitely great'') of various orders, appears implicitly from the first writings on differential Calculus, and for example in those of Fermat; it becomes fully conscious with Newton and Leibniz, with the theory of ``differences of higher order''; and no time is lost in observing that, in the simplest cases, the limit (or ``true value'') of an expression of the form \(f(x)/g(x)\) , at a point where f and g both tend to 0, is given by the Taylor expansion of these functions in a neighbourhood of the point under consideration (``l'Hôpital's rule'', due probably to Johann Bernoulli).

Apart from this elementary case, the main problem of ``asymptotic expansion'' which presents itself to mathematicians already at the end of the XVIIth century is the calculation, exact or approximate, of sums of the form \(\sum_{k=1}^{n}f(k)\) , when n is very large; such a calculation is in effect necessary as much for interpolation and the numerical evaluation of the sum of a series, as in the Calculus of probabilities, where the ``functions of large numbers'' such as \(n!\) or \(\binom{a}{n}\) play a preponderant role. Already Newton, in order to obtain approximate values of \(\sum_{k=1}^{n}\frac{1}{a+k}\) when n is large, indicates a method that reduces (in this particular case) to calculating the first terms of the Euler-Maclaurin formula ({[}262{]}, v. II, pp.~309-310). Towards the end of the century, Jakob Bernoulli, in the course of his research on the Calculus of probabilities, proposes determining the sums \(S_{k}(n)=\sum_{p=1}^{n}p^{k}\) , polynomials \(^{1}\) in n for which he discovers the general formation law (without giving a proof of it), introducing thus for the first time, in the expression for the coefficients of these polynomials, the numbers that bear his name, and the recurrence relation that enables their calculation ({[}19 b{]}, p.~97). In 1730, Stirling obtains an asymptotic expansion for \(\sum_{k=1}^{n}\log(x+ka)\) , n growing indefinitely, by a procedure for the calculation of the coefficients by recurrence.

From 1730 to 1745 is dated the decisive work of Euler on series and the relevant questions. Setting \(S(n) = \sum_{k=1}^{n} f(k)\) , he applies Taylor's formula to the function \(S(n)\) , which gives him

\[
f (n) = S (n) - S (n - 1) = \frac {d S}{d n} - \frac {1}{2 !} \frac {d ^ {2} S}{d n ^ {2}} + \frac {1}{3 !} \frac {d ^ {3} S}{d n ^ {3}} - \dots ,
\]

an equation that he ``inverts'' by the method of indeterminate coefficients, while looking for a solution of the form

\[
S (n) = \alpha \int f (n) d n + \beta f (n) + \gamma \frac {d f}{d n} + \delta \frac {d ^ {2} f}{d n ^ {2}} + \dots ;
\]

he thus obtains step by step

\[
S (n) = \int f (n) d n + \frac {f (n)}{2} + \frac {1}{12} \frac {d f}{d n} - \frac {1}{720} \frac {d ^ {3} f}{d n ^ {3}} + \frac {1}{30,240} \frac {d ^ {5} f}{d n ^ {5}} - \dots
\]

without at first being able to determine the formation law for the coefficients ({[}108 a{]}, (1), v. XIV, pp.~42-72 and 108-123). But around 1735, by analogy with the decomposition of a polynomial into factors of degree one, he does not hesitate to write the formula

\[
\begin{array}{c} 1 - \frac {\sin s}{\sin \alpha} = \\ \left(1 - \frac {s}{\alpha}\right) \left(1 - \frac {s}{\pi - \alpha}\right) \left(1 - \frac {s}{- \pi - \alpha}\right) \left(1 - \frac {s}{2 \pi + \alpha}\right) \left(1 - \frac {s}{- 2 \pi + \alpha}\right) \dots \end{array}
\]

and in putting the coefficients of the expansions of the two members of the whole series equal, he obtains in particular (for \(\alpha = \frac{\pi}{2}\) ) the famous expressions for the series \(\sum_{n=1}^{\infty} \frac{1}{n^{2k}}\) by means of the powers \(^2\) of \(\pi\) (loc. cit., pp.~73-86). A few years later, he realises at last that the coefficients of these powers of \(\pi\) are given by the same equations as those of his summation formula, and recognises their link with the numbers introduced by Bernoulli, and with the coefficients of the expansion in series of \(z/(e^z - 1)\) (loc. cit., pp.~407-462).

Independently of Euler, Maclaurin had arrived at about the same time at the same summing formula, by a somewhat less hazardous path, close to that followed today: he iterates in effect the ``Taylorian'' formula which expresses \(f(x)\) by means of the differences \(f^{(2k + 1)}(x + 1) - f^{(2k + 1)}(x)\) , a formula that he obtains by ``inverting'' the Taylor expansions of these differences by the method of indeterminate coefficients ({[}214{]}, v. II, pp.~672-675); he does not notice besides the formation law for the coefficients, discovered by Euler.

But Maclaurin, like Euler and all the mathematicians of his time, presents all his formulae as expansions in series, whose convergence is not even studied. It is not that the notion of convergent series was totally neglected at this time: it is known from Jakob Bernoulli that the harmonic series is divergent, and Euler had himself recorded this result while evaluating the sum of the first n terms of this series by means of his summing formula ({[}108 a{]}, (1), v. XIV, pp.~87-100 and 108-123); it is also Euler who remarks that the ratio of two consecutive Bernoulli numbers grows indefinitely, and that it follows that an entire series having these numbers as coefficients can not converge (loc. cit., p.~357). \(^{3}\) But the tendency to a formal calculation is the stronger, and the extraordinary intuition of Euler himself does not prevent him from lapsing sometimes into the absurd, when he writes for example \(0 = \sum_{n=-\infty}^{+\infty} x^{n}\) (loc. cit., p.~362). \(^{4}\)

We have said elsewhere (see p.~153) how the mathematicians of the beginning of the XIXth century, tired of this formalism without restraint and without basis, brought Analysis back into the paths of rigour. Once the notion of a convergent series is made precise, the necessity appeared for simple criteria providing a proof of the convergence of integrals and series by comparison with known integrals or series; Cauchy gives a certain number of criteria in his Analyse algébrique ({[}56 a{]}, (2), v. III), whereas Abel, in a posthumous memoir ({[}1{]}, v. II, pp.~197-205), obtains logarithmic criteria of convergence. Cauchy, on the other hand ({[}56 a{]}, (1), v. VIII, pp.~18-25), elucidates the paradox of series such as the Stirling series, obtained by the application of the Euler-Maclaurin formula (and often called ``semi-convergent series''): he shows that if (because of the remark of Euler on the Bernoulli numbers) the general term \(u_k(n)\) of such a series, for a fixed value of \(n\) , grows indefinitely with \(k\) , it remains none the less that, for a fixed value of \(k\) , the partial sum \(s_k(n) = \sum_{h=1}^{k} u_h(n)\) gives an asymptotic expansion (for \(n\) tending to \(+\infty\) ) of the function ``represented'' by the series, all the more accurate as \(k\) is larger.

In the majority of the calculations of classical Analysis, it is possible to obtain a general formation law for asymptotic expansions of a function, having an arbitrarily large number of terms; this fact contributed to creating a lasting confusion (at least in the language) between series and asymptotic expansions; so much so that H. Poincaré, when he takes the trouble, in 1886 ({[}251 a{]}, v. I, pp.~290-296), of codifying the elementary rules of asymptotic expansions (following the integral powers of \(1/x\) in the neighbourhood of \(+\infty\) ), uses still the vocabulary of the theory of series. It is only with the appearance of asymptotic expansions coming from the analytic theory of numbers that there was a clear distinction operating between the notion of asymptotic expansion and that of series, by reason of the fact that, in the majority of problems that this theory deals with, only a very small number of terms (most often only one) in the sought for expansion can be obtained explicitly.

These problems also made mathematicians familiar with the use of scales of comparison other than those of powers of the variable (real or integral). This extension goes back to the work of P. du Bois-Reymond {[}94 a and b{]} who, first, tackled systematically the problems of comparing functions in the neighbourhood of a point, and, in very original works, recognised the ``non-Archimedean'' character of the scales of comparison, at the same time as he studied in a general way the integration and differentiation of the comparison relations, and drew from this a large number of interesting consequences {[}94 b{]}. His proofs however lacked clarity and rigour, and it is to G. H. Hardy {[}147{]} that the correct presentation of the results of du Bois-Reymond is due: his main contribution consisted in recognising and proving the existence of a set of ``elementary functions'', the functions (H), where the usual operations of Analysis (notably differentiation) are applicable to comparison relations.

\chapter{THE GAMMA FUNCTION.}

The idea of ``interpolating'' a sequence \((u_{n})\) by the values of an integral dependant on a real parameter \(\lambda\) and equal to \(u_{n}\) for \(\lambda = n\) , goes back to Wallis (cf.~pp.~187 ff.). It is this idea that was the main guide for Euler when, in 1730 ({[}108 a{]}, (1), v. XIV, pp.~1-24), he has the idea of interpolating the sequence of factorials. He starts by remarking that \(n!\) is equal to the infinite product \(\prod_{k=1}^{\infty}\left(\frac{k+1}{k}\right)^{n}\frac{k}{k+n}\) , that this product is defined for every value of n (integral or not), and that in particular, for \(n = \frac{1}{2}\) it takes the value \(\frac{1}{2}\sqrt{\pi}\) from the formula of Wallis. The analogy of this result with those of Wallis leads him then to take up again the integral \(\int_{0}^{1} x^{e}(1-x)^{n} dx\) (n an integer, e arbitrary), already considered by this latter. Euler obtains its value \(\frac{n!}{(e+1)(e+2)\cdots(e+n)}\) by the binomial expansion; a change of variable then shows him that \(n!\) is the limit, for z tending to 0, of the integral \(\int_{0}^{1}\left(\frac{1-x^{s}}{z}\right)^{n} dx\) , whence the ``second Eulerian integral'' \(n! = \int_{0}^{1}\left(\log\frac{1}{x}\right)^{n} dx\) ; by the same method and the use of the formula of Wallis, he obtains the formula \(\int_{0}^{1}\sqrt{\log\frac{1}{x}} dx = \frac{1}{2}\sqrt{\pi}\) . In his later work, Euler returns frequently to these integrals; he discovers in this way the relation of the complements ({[}108 a{]}, (1), v. XV, p.~82 and v. XVII, p.~342), the formula \(B(p,q) = \Gamma(p)\Gamma(q)/\Gamma(p+q)\) ({[}108 a{]}, (1), v. XVII, p.~355), and the particular case of the Legendre-Gauss formula corresponding to x = 1 ({[}108 a{]}, (1), v. XIX, p.~483); all of it of course without worrying about questions of convergence.

Gauss continued the study of the \(\Gamma\) function in connection with his research on the hypergeometric function, of which the function \(\Gamma\) is a limit case ({[}124 a{]}, v. III, pp.~125-162); it is during this research that he obtains the general formula for multiplication (already noted by Legendre not long before for \(p = 2\) ). The subsequent work on \(\Gamma\) has mainly been concerned with the extension of this function to the complex domain. It is only recently that it has been noticed that the logarithmic convexity property characterised \(\Gamma(x)\) (in the real domain) up to a factor amongst all the solutions of the functional equation \(f(x + 1) = xf(x)\) ({[}26{]}, pp.~149-164); and Artin has shown {[}7 d{]} how all the classical results about \(\Gamma(x)\) can be simply linked to this property.

\chapter{FUNCTION SPACES.}

We know that the notion of an arbitrary function did not emerge much before the beginning of the XIXth century. All the more so the idea of studying in a general way sets of functions, and furnishing them with a topological structure, did not appear before Riemann (see p.~141) and does not begin to be put to use before the end of the XIXth century.

However, the notion of convergence of a sequence of numerical functions was used in a more or less conscious way since the beginnings of the infinitesimal Calculus. But it was only a matter of simple convergence and it could not be otherwise before the notions of convergent series and of continuous functions had been defined accurately by Bolzano and Cauchy. This latter did not see at the first attempt the distinction between simple convergence and uniform convergence, and believed it possible to prove that every convergent series of continuous functions has as sum a continuous function ({[}56 a{]}, (2), vol.~III, p.~120). The error was almost immediately uncovered by Abel, who proved at the same time that every entire series is continuous in the interior of its convergence interval, by means of the argument that has become a classic which uses essentially, in this particular case, the idea of uniform convergence ({[}1{]}, v. I, pp.~223-224). It only remained to disengage this latter in a general way, which was done independently by Stokes and Seidel in 1847-48, and by Cauchy himself in 1853 ({[}56 a{]}, (1), vol.~XII, p.~30). \(^{1}\)

Under the influence of Weierstrass and of Riemann, the systematic study of the notion of uniform convergence and of connected questions is developed in the last third of the XIXth century by the German school (Hankel, du Bois-Reymond) and especially the Italian school: Dini and Arzelà specify the necessary conditions for the limit of a sequence of continuous functions to be continuous, whereas Ascoli introduces the fundamental notion of equicontinuity and proves the theorem characterising compact sets of continuous functions {[}10{]} (a theorem popularised later by Montel in his theory of ``normal families'', which are none other than relatively compact sets of analytic functions).

Weierstrass himself discovered on the other hand ({[}329 a{]}, v. III, pp.~1-37) the possibility of approximating uniformly by polynomials a continuous numerical function of one or several real variables on a bounded set, a result that aroused immediately a lively interest, and led to numerous ``quantitative'' studies that remain outside the point of view we follow at present.

The modern contribution to these questions consisted above all in giving them the whole range of which they are capable, by bringing them into play for functions whose set of definition and set of values are no longer restricted to R or to spaces of finite dimension, and in placing them thus in their natural framework, with the help of general topological concepts. In particular, the theorem of Weierstrass, which had already shown itself to be a tool of the first order in classical Analysis, could, in these latter years, be extended to much more general cases by M. H. Stone: developing an idea introduced by H. Lebesgue (in a proof of the theorem of Weierstrass), he highlighted the important role played in the approximation of continuous numerical functions by lattices (approximation by ``lattice polynomials''), and showed on the other hand how the generalised theorem of Weierstrass implies immediately a whole series of analogous approximation theorems, that are thus grouped in a much more coherent way {[}301 c{]}.

\chapter{TOPOLOGICAL VECTOR SPACES.}

The general theory of topological vector spaces was founded during the period which goes from 1920 to 1930 approximately. But it had been prepared for a long time before by the study of numerous problems of Functional Analysis; and its history cannot be retraced without indicating, at least summarily, how the study of these problems gradually brought mathematicians (especially from the beginning of the XXth century) to become aware of the common ancestry of the questions under consideration, and of the possibility of formulating them in a much more general way and of applying to them uniform solution procedures.

It can be said that the analogies between Algebra and Analysis, and the idea of considering functional equations (that is to say where the unknown is a function) as ``limit cases'' of algebraic equations, go back to the beginnings of infinitesimal Calculus, which in a kind of way answers this need to generalise ``from the finite to the infinite''. But the direct algebraic ancestor of the infinitesimal Calculus is the calculus of finite differences (cf.~pp.~185 ff.), and not the solution of general linear systems; and it is not before the middle of the XVIIIth century that the first analogies between this latter and the problems of the differential Calculus are manifested, in connection with the equation of vibrating cords. We will not go into the detail of the history of this problem here; but we must pick up the appearance of two fundamental ideas, which are found again constantly in what follows, and which both appear to be due to D. Bernoulli. The first consists in considering the oscillation of the cord as a ``limit case'' of the oscillation of a system of n point masses, as n tends to infinity; it is known that, for n finite, this problem would a little bit later give the first example of the search for eigenvalues of a linear transformation (cf.~p.~88); there correspond to these numbers, in the ``passage to the limit'' being considered, the frequencies of the ``proper oscillations'' of the cord, observed experimentally since long before, and whose theoretic existence had been established (notably by Taylor) at the beginning of the century. This formal analogy, although fairly rarely mentioned in what followed ({[}302 b{]}, p.~390), seems never to have been lost sight of during the XIXth century; but, as we will see later, it did not acquire its full importance until about 1890-1900.

The other idea of D. Bernoulli (perhaps inspired by experimental facts) is the ``principle of superposition'', according to which the most general oscillation of the cord must be capable of being ``decomposed'' into the superposition of ``proper oscillations'', which, mathematically speaking, means that the general solution of the equation of vibrating cords must be capable of being expanded as a series \(\sum_{n} c_{n} \phi_{n}(x, t)\) , where the \(\phi_{n}(x, t)\) represent the proper oscillations. It is known that this principle was to let loose a long quarrel on the possibility of expanding an ``arbitrary'' function as a trigonometric series, a quarrel that was only settled by the work of Fourier and of Dirichlet in the first third of the XIXth century. But even before this result was reached, other examples had been met of expansion in series of ``orthogonal'' \(^{1}\) functions: spherical functions and Legendre polynomials, as well as different systems of the form \((e^{i\lambda_{n}x})\) , where the \(\lambda_{n}\) are no longer multiples of the same number, and which were introduced from the XVIIIth century in oscillation problems, as well as by Fourier and Poisson during their researches on the theory of heat. Around 1830, all the phenomena observed in these various special cases are systematised, by Sturm {[}302 a and b{]} and Liouville {[}204 a and b{]} in a general theory of oscillations, for functions of one variable: they consider the differential equation

\[
\frac {d}{d x} \left(p (x) \frac {d y}{d x}\right) + \lambda \rho (x) y = 0 (p (x) > 0, \rho (x) > 0)\tag{21.1}
\]

with boundary conditions

\[
\left\{ \begin{array}{l l} y ^ {\prime} (a) - h _ {1} y (a) & = 0 \\ y ^ {\prime} (b) + h _ {2} y (b) & = 0 \end{array} \right. (h _ {1} \geq 0, h _ {2} \geq 0, a <   b)\tag{21.2}
\]

and prove the following fundamental results:

\begin{enumerate}
\def\labelenumi{\arabic{enumi})}
\item
  the problem only has a solution \(\neq 0\) when \(\lambda\) takes one of the values from a sequence \((\lambda_{n})\) of numbers \textgreater{} 0, tending to \(+\infty\) ;
\item
  for each \(\lambda_{n}\) , the solutions are multiples of the same function \(v_{n}\) , that can be assumed to be ``normalised'' by the condition \(\int_{a}^{b} \rho v_{n}^{2} dx = 1\) , and we have \(\int_{a}^{b} \rho v_{m} v_{n} dx = 0\) for \(m \neq n\) ;
\item
  every function \(f\) , twice differentiable in \([a, b]\) , and satisfying the boundary conditions (21.2), is expandable as a uniformly convergent series \(f(x) = \sum_{n} c_{n} v_{n}(x)\) , where \(c_{n} = \int_{a}^{b} \rho f v_{n} dx\) ;
\item
  we have the equality \(\int_{a}^{b}\rho f^{2}dx=\sum_{n}c_{n}^{2}\) (already proved by Parseval in 1799 --- in a purely formal way, besides --- for the system of trigonometric functions, and from which flows immediately the ``Bessel inequality'' stated by this latter (still for trigonometric series) in 1828).
\end{enumerate}

Half a century later, these properties are completed by the work of Gram {[}133{]} who, following the researches of Tchebychef, brings to light the relation between the expansions in series of orthogonal functions and the problem of the ``best quadratic approximation'' (issuing directly from the ``method of least squares'' of Gauss, in error theory): this latter consists, given a finite sequence of functions \((\psi_{i})_{1\leq i\leq n}\) , of finding, for a function f, the linear combination \(\sum_{i}a_{i}\psi_{i}\) for which the integral \(\int_{a}^{b}\rho(f-\sum_{i}a_{i}\psi_{i})^{2}dx\) attains its minimum. In principle it is a case there only of a basic linear algebra problem, but Gram solves it in an original way, by applying to the \(\psi_{i}\) the process of ``orthonormalization'' generally known by the name of Erhard Schmidt. Going on next to the case of an infinite orthogonal system \((\phi_{n})\) , he sets himself the question of knowing when the ``best quadratic approximation'' \(\mu_{n}\) of a function f by linear combinations of the first n functions of the sequence, tends to 0 when n tends to infinity; \(^{2}\) he is thus led to define the notion of a complete orthonormal system, and recognises that this property is equivalent to the non-existence of functions \(\neq0\) orthogonal to all the \(\phi_{n}\) . He even seeks to elucidate the concept of ``mean quadratic convergence'', but, before the introduction of the fundamental notions of measure theory, he could hardly obtain any very specific results in this direction.

In the second half of the XIXth century, the main effort of the analysts is mainly towards the extension of the Sturm-Liouville theory to functions of several variables, to which led notably the study of partial differential equations of elliptic type from mathematical Physics, and the boundary problems that are naturally associated with them. Interest is concentrated principally on the ``vibrating membranes'' equation

\[
L _ {\lambda} (u) \equiv \Delta u + \lambda u = 0\tag{21.3}
\]

where the solutions are sought that are zero on the contour for a fairly regular domain \(G\) ; and it is only little by little that the considerable analytic difficulties presented by this problem were overcome, a problem to which it was not thought possible to apply the methods that had succeeded for functions of a single variable. Let us remind ourselves of the main stages towards the solution: the introduction of the ``Green's function'' of \(G\) , whose existence was proved by Schwarz; the proof, also due to Schwarz, of the existence of the smallest eigenvalue; finally, in 1894, in a famous memoir ({[}251 a{]}, v. IX, pp.~123-196), H. Poincaré succeeded in proving the existence and essential properties of all the eigenvalues, by considering, for a given ``second member'' \(f\) , the solution \(u_{\lambda}\) of the equation \(L_{\lambda}(u) = f\) which is zero on the contour, and in proving, by a clever generalisation of the method of Schwarz, that \(u_{\lambda}\) is a meromorphic function of the complex variable \(\lambda\) , having only real simple poles \(\lambda_n\) , which are precisely, the sought for eigenvalues.

This research is closely linked to the beginnings of the theory of linear integral equations, which was to contribute no doubt most to the coming of modern ideas. We will limit ourselves here to giving a few short indications on the development of this theory. This type of functional equation, having first appeared sporadically in the first half of the XIXth century (Abel, Liouville), had acquired importance since Beer and C. Neumann had reduced the solution of ``Dirichlet's problem'' for a fairly regular domain G, to the solution of an ``integral equation of the second type''

\[
u (x) + \int_ {a} ^ {b} K (x, y) u (y) d y = f (x)\tag{21.4}
\]

for an unknown function u; an equation that C. Neumann had succeeded in solving by means of a procedure of ``successive approximations'' in 1877. Driven no doubt as much by the algebraic analogies already mentioned as by the results that he had just obtained on the vibrating membranes equation, H. Poincaré, in 1896 ({[}251 a{]}, v. IX, pp.~202-272), had the idea of introducing a variable parameter \(\lambda\) in front of the integral in the preceding equation, and affirms that, as for the vibrating membranes equation, the solution is then a meromorphic function of \(\lambda\) ; but he does not succeed in proving this result, which was only established (for a continuous ``kernel'' K and a finite interval \([a, b]\) ) by I. Fredholm four years later {[}116{]}. This latter, more consciously perhaps even than his predecessors, lets himself be completely guided by the analogy of (21.4) with the linear system

\[
\sum_ {q = 1} ^ {n} \left(\delta_ {p q} + \frac {1}{n} a _ {p q}\right) x _ {q} = b _ {p} (1 \leq p \leq n)\tag{21.5}
\]

in order to obtain the solution of (21.4) as a quotient of two expressions constructed on the model of determinants that occur in the formulae of Cramer. It was not besides a new idea: from the beginning of the XIXth century, the method of ``indeterminate coefficients'' (consisting of obtaining an unknown function assumed expandable in series \(\sum_{n} c_{n} \phi_{n}\) , where the \(\phi_{n}\) are known functions, by calculating the coefficients \(c_{n}\) ) had led to ``linear systems in an infinity of unknowns''

\[
\sum_ {j = 1} ^ {\infty} a _ {i j} x _ {j} = b _ {i} (i = 1, 2, \dots).\tag{21.6}
\]

Fourier, who comes across such a system, ``solves'' it still like a mathematician of the XVIIIth century: he suppresses all the terms having an index i or j greater than n, solves explicitly the finite system thus obtained, by means of Cramer's rule, then ``passes to the limit'' by making n tend to +∞ in the solution! When later it was no longer sufficient to play such lazy games of pass the parcel, it is still by means of the theory of determinants that the first attempts to attack the problem were made; starting in 1886 (following the work of Hill), H. Poincaré, then H. von Koch, had built a theory of ``infinite determinants'' which allows the solution of certain types of systems (21.6) following the classical model; and if these results were not directly applicable to the problem Fredholm aimed at, at least it is certain that the theory of von Koch, in particular, was of use as a model for the formation of his ``determinants''.

It is at this moment that Hilbert enters the scene and gives a new impulsion to the theory {[}163 b{]}. He starts by completing the work of Fredholm in achieving effectively the passage to the limit that leads from the solution of (21.5) to that of (21.4); but he adds to it immediately the corresponding passage to the limit for the theory of real quadratic forms, to which the types of integral equations with symmetric kernel (that is to say such that \(K(y, x) = K(x, y)\) ) were leading naturally, by far the most frequent in mathematical Physics. He reaches thus the fundamental formula that generalises directly the reduction of a quadratic form to its axes

\[
\int_ {a} ^ {b} \int_ {a} ^ {b} K (s, t) x (s) x (t) d s d t = \sum_ {n = 1} ^ {\infty} \frac {1}{\lambda_ {n}} \left(\int_ {a} ^ {b} \phi_ {n} (s) x (s) d s\right) ^ {2},\tag{21.7}
\]

the \(\lambda_{n}\) being the eigenvalues (necessarily real) of the kernel \(K\) , the \(\phi_{n}\) forming the orthogonal system of corresponding eigenfunctions, and the second member of the formula (21.7) being a convergent series for \(\int_{a}^{b} x^{2}(s) dx \leq 1\) . He shows also how every function ``representable'' in the form \(f(x) = \int_{a}^{b} K(x, y) g(y) dy\) allows the ``expansion'' \(\sum_{n=1}^{\infty} \phi_{n} \int_{a}^{b} \phi_{n}(y) f(y) dy\) , and carrying on the analogy with the classical theory of quadratic forms, he indicates a procedure for the determination of the \(\lambda_{n}\) by a variational method, which is none other than the extension of the well-known extremal properties of the axes of a quadric ({[}163 b{]}, pp.~1-38).

These first results of Hilbert were almost immediately taken up by E. Schmidt, in a simpler and more general form, avoiding the introduction of the ``Fredholm determinants'' as well as the passage from finite to infinite, and already very near an abstract statement, the fundamental properties of linearity and of positivity of the integral being visibly used only in the proofs {[}274 a{]}. But already Hilbert had reached conceptions which are much more general still. All the preceding work brought out the importance of functions that are square integrable, and the Parseval formula established a tight link between these functions and the sequences \((c_{n})\) such that \(\sum_{n} c_{n}^{2} < +\infty\) . It is no doubt this idea that guides Hilbert in his memoirs of 1906 ({[}163 b{]}, chap.~XI-XIII) where, taking up the old method of ``indeterminate coefficients'', he shows that the solution of the integral equation (21.4) is equivalent to a system of an infinity of linear equations

\[
x _ {p} + \sum_ {q = 1} ^ {\infty} k _ {p q} x _ {q} = b _ {p} (p = 1, 2, \dots)\tag{21.8}
\]

for the ``Fourier coefficients'' \(x_{p} = \int_{a}^{b} u(t) \omega_{p}(t) dt\) of the unknown function u with respect to a given complete orthonormal system \((\omega_{n})\) (with \(b_{p} = \int_{a}^{b} f(t) \omega_{p}(t) dt\) and \(k_{pq} = \int_{a}^{b} \int_{a}^{b} K(s, t) \omega_{p}(s) \omega_{q}(t) ds dt\) . As well, the only solutions of (21.8) to be considered from this point of view are those for which \(\sum_{n} x_{n}^{2} < +\infty\) ; also it is to this type of solution that Hilbert systematically limits himself; but over against that he relaxes the conditions imposed on the ``infinite matrix'' ( \(k_{pq}\) ) (which, in (21.8), is such that \(\sum_{p,q} k_{pq}^{2} < +\infty\) ). From that moment, it is clear that the ``Hilbert space'' of sequences \(x = (x_{n})\) of real numbers such that \(\sum_{n} x_{n}^{2} < +\infty\) , although not explicitly introduced, underlies all the theory, and appears as a ``passage to the limit'' starting from Euclidean space of finite dimension. Over and above that, what is particularly important for subsequent developments, Hilbert is led to introduce in this space, not only one, but two distinct notions of convergence (corresponding to what has been called the weak topology and the strong topology \(^{3}\) ), as well as a ``principle of choice'' which is none other than the property of weak compactness of the unit ball. The new linear algebra that he develops in connection with the solution of the systems (21.8) rests entirely on these topological notions: linear maps, linear forms and bilinear forms (associated with linear maps) are classified and studied according to their ``continuity'' \(^{4}\) properties. In particular, Hilbert discovers that the success of Fredholm's method rests on the notion of ``complete continuity'', that he brings out by formulating it for bilinear forms \(^{5}\) and studying it deeply; we can not give here more details about this important notion, nor on the admirable and deep work where Hilbert inaugurates the spectral theory of symmetric bilinear forms (bounded or not).

The language of Hilbert still remains classical, and, throughout the ``Grundzüge'', he never ceases to have in view the applications of the theory, of which he develops numerous examples (taking up about half the volume). The following generation will already adopt a more abstract point of view. Under the influence of ideas of Fréchet and of F. Riesz on general topology (see p.~142), E. Schmidt {[}274 b{]} and Fréchet himself introduce deliberately, in 1907-1908, the language of Euclidean geometry into ``Hilbert space'' (real or complex); it is in these works that is found the first mention of the norm (with the actual notation \(||x||\) ), the triangle inequality that it satisfies, the fact that Hilbert space is ``separable'' and complete; moreover, E. Schmidt proves the existence of the orthogonal projection on a closed linear variety, which allows him to give a simpler and more general form to Hilbert's theory of linear systems. In 1907 also, Fréchet and F. Riesz remark that the space of square summable series has a fully analogous ``geometry''; an analogy that is explained perfectly when, some months later, F. Riesz and E. Fischer prove that this space is complete and isomorphic to ``Hilbert space'', at the same time bringing out in a startling way the value of the tool newly created by Lebesgue. From that moment the essential points of the theory of Hilbert spaces can be considered as acquired; among more recent progress, mention should most notably be made of the axiomatic presentation of the theory given around 1930 by M. H. Stone and J. von Neumann, as well as the casting off of the restrictions of ``separability'', which is effected around 1934, in the works of Rellich, Löwig, and F. Riesz {[}260 g{]}.

However other streams of ideas were coming, in the first years of the XXth century, to reinforce the tendency that was leading to the theory of normed spaces. The general idea of a ``functional'' (that is to say a function with numerical values defined on a set whose elements are themselves numerical functions of one or several real variables) was brought out in the last decades of the XIXth century, linked to the calculus of variations, on the one hand, the theory of integral equations on the other. But if it is mainly to the Italian school, around Pincherle and especially Volterra, that is due the bringing to light of this notion, as well as the more general idea of ``operator'', the work of this school remained often of a passably formal nature and attached to problems of a particular type, lacking a sufficiently deep analysis of the underlying topological concepts. In 1903, Hadamard inaugurates the modern theory of ``topological'' duality, by seeking the most general continuous linear ``functionals'' on the space \(\mathcal{C}(I)\) of numerical functions on a compact interval I (a space furnished with a topology of uniform convergence), and in characterising them as limits of sequences of integrals \(x \mapsto \int_{I} k_{n}(t)x(t)dt\) . In 1907, Fréchet and F. Riesz show in the same way that, on a Hilbert space, continuous linear forms are the ``bounded'' forms introduced by Hilbert; then, in 1909, F. Riesz puts into definitive form Hadamard's theorem, by expressing all continuous linear functionals on \(\mathcal{C}(I)\) by a Stieltjes integral, a theorem that would later become the point of departure for the modern theory of integration (see p.~226).

The following year, it is still F. Riesz {[}260 c{]} who makes new and important progress in the theory by introducing and studying (copied from the theory of Hilbert space) the spaces \(L^{p}(I)\) of functions that are p-power summable in an interval I (for an exponent p such that \(1 < p < +\infty\) ), a study that he follows three years later {[}260 e{]} with analogous work on the spaces of sequences \(L^{p}(\mathbf{N})\) ; these researches, as we will see later, were greatly to contribute to clarifying the ideas on duality, because of the fact that here were met for the first time two spaces which were dual and not naturally isomorphic. \(^{6}\)

From that moment, F. Riesz was thinking about an axiomatic study encompassing all these results ({[}260 c{]}, p.~452); and it seems that only the scruples of an analyst concerned not to distance himself too much from classical mathematics kept him from writing his famous memoir of 1918 on Fredholm theory in this form {[}260 f{]}. He considers there in principle the space \(\mathcal{C}(I)\) of functions continuous on a compact interval; but, after having defined the norm of this space, and having remarked that \(\mathcal{C}(I)\) , together with this norm, is complete, he never uses again in his arguments any thing other than the axioms of normed complete spaces. \(^{7}\) Without going into a detailed examination of this work here, let us mention that it is there that is to be found for the first time defined in a general way the notion of a completely continuous linear map (by the property of transforming a neighbourhood into a relatively compact set); \(^{8}\) and, by a chef-d'oeuvre of axiomatic analysis, the whole of Fredholm's theory (in its qualitative aspect) is reduced to a single fundamental theorem, namely that every normed locally compact space is finite dimensional.

The general definition of normed spaces was given in 1920-22 by S. Banach, H. Hahn and E. Helly (this latter only considering spaces of sequences of real or complex numbers). In the ten years that followed, the theory of spaces is developed mainly around two questions of fundamental importance for applications: the theory of duality and the theorems that are attached to the notion of Baire ``category''.

We have seen that the idea of duality (in the topological sense) goes back to the beginning of the XXth century; it underlies Hilbert's theory and occupies a central place in the work of F. Riesz. This latter observes, for example, already in 1911 ({[}260 d{]}, pp.~41-42), that the relation \(|f(x)| \leq M||x||\) (taken as the definition of ``bounded'' linear functionals in Hilbert space) is equivalent to the continuity of f when one is in the space \(\mathcal{C}(I)\) , and that by an argument of fully general character. In connection with the characterisation of continuous linear functionals on \(\mathcal{C}(I)\) , he remarks also that the condition for a set A to be dense in \(\mathcal{C}(I)\) is that there does not exist any Stieltjes measure \(\mu \neq 0\) on I that is ``orthogonal'' to every function of A (generalising thus the Gram condition for complete orthonormal systems); he notes finally, in the same work, that the dual of the space \(L^{\infty}\) is ``bigger'' than the space of Stieltjes measures ({[}260 d{]}, p.~62).

On the other hand, in his work on the spaces \(L^{p}(I)\) and \(L^{p}(\mathbf{N})\) , F. Riesz succeeds in modifying the method for the solution of linear systems in Hilbert space, given by E. Schmidt {[}274 b{]}, in such a way as to make it applicable to more general cases. The idea of E. Schmidt consisted in determining an ``extremal'' solution of (21.6) by seeking the point of the closed linear variety represented by the equations (21.6), whose distance to the origin is minimal. By using the same idea, F. Riesz shows that a necessary and sufficient condition that there exists a function \(x \in L^{p}(a, b)\) satisfying the equations

\[
\int_ {a} ^ {b} \alpha_ {i} (t) x (t) d t = b _ {i} (i = 1, 2, \dots)\tag{21.9}
\]

(where the \(\alpha_{i}\) belong to \(L^{q}\) (with \(\frac{1}{p} +\frac{1}{q} = 1\) ), and such that as well \(\int_{a}^{b}|x(t)|^{p}dt\leq M^{p}\) , is that, for all finite sequences \((\lambda_i)_{1\leq i\leq n}\) of real numbers, one has

\[
\left| \sum_ {i = 1} ^ {n} \lambda_ {i} b _ {i} \right| \leq M. \left(\int_ {a} ^ {b} \sum_ {i = 1} ^ {n} | \lambda_ {i} \alpha_ {i} (t) | ^ {q} d t\right) ^ {1 / q}.\tag{21.10}
\]

In 1911 {[}260 d{]}, he treats in an analogous way the ``problem of generalised moments'', consisting of solving the system

\[
\int_ {a} ^ {b} \alpha_ {i} (t) d \xi (t) = b _ {i} (i = 1, 2, \dots)\tag{21.11}
\]

where the \(\alpha_{i}\) are continuous and the unknown is a Stieltjes measure \(\xi;^{9}\) and it is visible here that the problem can be stated by saying that it is a case of determining a continuous linear functional on \(\mathcal{C}(I)\) by its values on a sequence of given points in this space. It is also in this form that Helly treats the problem in 1912, --- obtaining the conditions of F. Riesz by a fairly different method and of wider applicability \(^{10}\) --- and that he takes it up again in 1921, under much more general conditions. Introducing the notion of norm (on spaces of sequences) as we have seen above, he remarks that this notion generalises that of the ``gauge'' of a convex body in space of n dimensions, used by Minkowski in his famous work on the ``geometry of numbers'' {[}221 a and b{]}. In the course of this work, Minkowski had also defined (in \(R^{n}\) ) the notions of support hyperplane and of ``support function'' {[}221 b{]}, and proved the existence of a support hyperplane at all boundary points of a convex body ({[}221 a{]}, pp.~33-35). Helly extends these notions to a space of sequences E, furnished with an arbitrary norm; he establishes a duality between E and the space \(E'\) of sequences \(u = (u_n)\) such that, for all \(x = (x_n) \in E\) , the series \((u_n x_n)\) should be convergent; \(\langle u, x \rangle\) designating the sum of this series, he defines in \(E'\) a norm by the formula \(\sup_{x \neq 0} |\langle u, x \rangle| / ||x||\) , which gives the support function in finite dimensional spaces. \(^{11}\) The solution of a system (21.6) in E, where each of the sequences \(u_i = (a_{ij})_{j \geq 1}\) is assumed to belong to \(E'\) , reduces, as Helly shows then, to solving successively the following two problems: \(1^\circ\) find a continuous linear form L on the normed space \(E'\) , such that \(L(u_i) = b_i\) for every index i, which, as he indicates, leads to conditions of type (21.10); \(2^\circ\) seek if such a linear form can be written \(u \mapsto \langle u, x \rangle\) for an \(x \in E\) . This last problem, as Helly observes, does not necessarily have a solution even when L exists, and he limits himself to giving some sufficient conditions that imply the existence of the solution \(x \in E\) in certain particular cases {[}156{]}.

These ideas acquire their definitive form in 1927, in a fundamental memoir of H. Hahn {[}142{]} whose results are found again (independently) by S. Banach two years later {[}14 c{]}. The procedure of Minkowski-Helly is applied by Hahn to an arbitrary normed space, and therefore gives to the dual the structure of a (complete) normed space; which allows Hahn immediately to consider successive duals of a normed space, and to state in a general way the problem of reflexive spaces, foreseen by Helly. But above all the crowning problem of the extension of a continuous linear functional while preserving its norm is definitively solved by Hahn in a completely general way, by a transfinite induction argument on the dimension --- thus giving one of the first examples of an important application of the axiom of choice to Functional Analysis. \(^{12}\) To these results, Banach adds a thorough study of the relations between a continuous linear map and its transpose, extending to general normed spaces results known until then only for the spaces \(L^{p}\) {[}260 c{]}, by means of a deep theorem on the weakly closed subsets of a dual; these results are expressed besides in a more striking way by using the notion of quotient space of a normed space, introduced a few years later by Hausdorff and Banach himself. Finally, it is again Banach who discovers the link between the weak compactness of the unit ball (observed in many particular cases, as we have pointed out above) and reflexivity, at least for spaces of countable type ({[}14 a{]}, p.~189). The theory of the duality of normed spaces can, from this moment, be considered as settled in its major aspects.

At the same time some theorems of paradoxical aspect were clarified, the first examples of which go back to the neighbourhood of 1910. Hellinger and Toeplitz had in effect proved in substance that year, that a sequence of bounded bilinear forms \(B_n(x,y)\) on a Hilbert space, whose values \(B_n(a,b)\) for every given pair \((a,b)\) are bounded (by a number dependant a priori on \(a\) and \(b\) ), is in fact uniformly bounded on every ball. Their proof proceeds by contradiction and consists in constructing a particular pair \((a,b)\) violating the hypothesis, by an induction method known since by the name of ``method of the sliding hump'', and which is still in use in many analogous questions. Already in 1905, Lebesgue had besides used an analogous procedure to prove the existence of continuous functions whose Fourier series diverges at certain points; and, the same year as Hellinger and Toeplitz, he applied the same method to prove that a weakly convergent series in \(L^1\) is bounded in the norm. \(^{13}\) These examples are multiplied in the years that follow, but without the introduction of any new ideas until 1927, the date when Banach and Steinhaus (with the partial collaboration of S. Saks) link these phenomena to the notion of thin sets and to the theorem of Baire on complete metric spaces, obtaining a general statement that covers all the previous particular cases {[}15{]}. The study of questions of ``category'' in normed complete spaces besides leads Banach, at the same time, to numerous other results on continuous linear maps; the most remarkable and no doubt the deepest is the theorem of the ``closed graph'' which, like the Banach-Steinhaus theorem, showed itself to be a tool of the first order in modern Functional Analysis {[}14 c{]}.

The publication of the treatise of Banach on ``Linear Operators'' {[}14 a{]} marks, one could say, the beginning of the adult age for the theory of normed spaces. All the results of which we have been speaking, as well as many others, are set out in this volume, in a way that is still a bit disorganised, but accompanied with many striking examples taken from varied domains of Analysis, and which seemed to forecast a brilliant future for the theory. As it happened, the work had considerable success, and one of its most immediate effects was the quasi-universal adoption of the language and notations used by Banach. But, in spite of a large amount of research undertaken for 40 years on Banach spaces ({[}203 bis{]}), if one makes an exception of the theory of Banach algebras and its applications to commutative and non-commutative harmonic analysis, the almost total absence of new applications of the theory to the big problems of classical Analysis has somewhat disappointed the hopes based upon it.

It is rather in the sense of an expanding and a more thrusting axiomatic analysis of the concepts relative to normed spaces that the most fruitful developments have occurred. Although the functional spaces met since the beginning of the XXth century arose for the most part furnished with a ``natural'' norm, it had not happened without some exceptions being noticed. Around 1910, E. H. Moore had proposed generalising the notion of uniform convergence by replacing it with the notion of ``relative uniform convergence'', where a neighbourhood of 0 is made up of the functions f satisfying a relation \(|f(t)| \leq \epsilon g(t)\) , g being a function everywhere \textgreater{} 0, which could vary with the neighbourhood. It had on the other hand been observed, before 1930, that notions such as simple convergence, convergence in measure for measurable functions, or compact convergence for entire functions, are not capable of being defined by means of a norm; and in 1926, Fréchet had noted that vector spaces of this nature can be metrisable and complete. But the theory of these more general spaces was only to develop in a fruitful way in combination with the idea of convexity. This latter (that we saw appearing with Helly) was an object of study for Banach and his pupils, who recognised the possibility of interpreting thus in a more geometric way numerous statements of the theory of normed spaces, preparing the way for the general definition of locally convex spaces, given by J. von Neumann in 1935. The theory of these spaces, and notably the questions dealing with duality, have been especially developed during the last ten years. What must be noted in this context, is on the one hand, the progress in simplicity and in generality made possible by the setting up of the fundamental notions of general Topology, carried out between 1930 and 1940; on the other hand, the importance taken by the notion of bounded set, introduced by Kolmogoroff and von Neumann in 1935, and whose fundamental role in duality theory has been brought to light by the work of Mackey {[}213 a and b{]} and of Grothendieck {[}138 b and c{]}. Finally and especially, it is certain that the main impulsion that motivated this research came from new possibilities for applications to Analysis, in domains where Banach theory was inoperative: the theory of sequence spaces must be mentioned in this context, developed by Köthe, Toeplitz and their pupils since 1934 in a series of memoirs {[}184{]}, the recent setting up of the theory of ``analytic functionals'' of Fantappié, and above all the theory of distributions of L. Schwartz {[}280{]}, where the modern theory of locally convex spaces found a field of applications that is without doubt a long way from being exhausted.

\chapter{INTEGRATION IN LOCALLY COMPACT SPACES.}

The development of the modern notion of integral is tightly linked to the evolution of the idea of function, and to a deep study of the numerical functions of real variables, which has been taking place since the beginning of the XIXth century. It is known that Euler already conceived of the notion of function in a fairly general way, since for him a given ``arbitrary'' curve intersected in a single point by every parallel to the axis Oy defined a function \(y = f(x)\) (cf.~p.~196); but, like the majority of his contemporaries, he refused to allow that such functions could be expressed ``analytically''. This point of view was to change hardly at all until the work of Fourier; but the discovery, by this latter, of the possibility of representing discontinuous functions as sums of trigonometric series, \(^{1}\) was to exercise a decisive influence on subsequent generations. Truth to say, the proofs of Fourier totally lacked in rigour, and their domain of validity did not appear clearly; however the integral formulae

\[
a _ {n} = \frac {1}{\pi} \int_ {- \pi} ^ {+ \pi} \phi (x) \cos n x d x, b _ {n} = \frac {1}{\pi} \int_ {- \pi} ^ {+ \pi} \phi (x) \sin n x d x (n \geq 1)\tag{22.1}
\]

giving the coefficients in the expansion of \(\phi\) as a Fourier series, had an intuitively obvious meaning as soon as \(\phi\) is assumed to be continuous and piecewise monotonic. \(^{2}\) Also it is at first to these functions that Dirichlet limits himself, in the famous memoir ({[}92{]}, v. I, pp.~117-132) where he established the convergence of Fourier series; but already, at the end of his work, he is concerned with the extension of his results to wider classes of functions. It is known that it is on this occasion that Dirichlet, making the ideas of Fourier precise, defines the general notion of function such as we understand it today; the first point to clarify was naturally to know in which cases it was still possible to attach a meaning to the formulae (22.1). ``When the solutions for the continuity {[}of \(\phi\) {]} are infinite in number \ldots{}'' says Dirichlet (loc. cit., pp.~131-132), ``it is necessary that then the function \(\phi(x)\) should be such that, if one denotes by \(a\) and \(b\) two arbitrary quantities between \(-\pi\) and \(+\pi\) , it is always possible to put between \(a\) and \(b\) other quantities \(r\) and \(s\) close enough so that the function remains continuous in the interval \(r\) to \(s\) . The necessity for this restriction will be felt by considering that the different terms of the series {[}of Fourier{]} are definite integrals and in going back to the fundamental notion of integrals. It will be seen then that the integral of a function does not mean anything except in as much as the function satisfies the condition previously stated.''

In modern terms, Dirichlet seems to believe that integrability is equivalent to the fact that the points of discontinuity make up a sparse set; he points out besides, a few lines later, the famous example of the function equal to c for x rational, and to a different value d for x irrational, and affirms that this function ``could not be substituted'' in the integral. He announced besides subsequent work on this subject, but this work was never published, \(^{3}\) and for 25 years, no one seems to have sought to advance along this path, perhaps because the consideration of such ``pathological'' functions seemed at the time completely lacking in interest; in any case, when Riemann, in 1854 ({[}259 a{]}, pp.~227-264), takes up the question (always in connection with trigonometric series \(^{4}\) ) he experiences the need to justify his work: ``Whatever our ignorance concerning the way in which the forces and the states of matter vary with time and place in the infinitely small, we can nonetheless take it as certain that the functions to which the researches of Dirichlet do not apply, do not occur in natural phenomena. However'', he continues, ``it appears that these cases not treated by Dirichlet deserve attention for two reasons. Firstly, as Dirichlet himself remarks at the end of his work, this subject is in a very close relation to the principles of the infinitesimal calculus, and can be of use to bring more clarity and certainty to these principles. From this point of view, his study has an immediate interest. In the second place, the application of Fourier series is not limited to research in physics; they are at present applied also with success in a domain of pure mathematics, number theory, and there it seems that they are precisely the functions whose expansion in trigonometric series has not been studied by Dirichlet, that are of importance'' ({[}259 a{]}, pp.~237-238).

The idea of Riemann is to start from the approximation procedure for the integral, brought back to respectability by Cauchy, and to determine when the ``Riemann sums'' of a function f, in a bounded interval \([a, b]\) , tend to a limit (the maximal length of the intervals in the subdivision tending to 0); a problem of which he obtains without trouble the solution in the following form: for every \(\alpha > 0\) , there is a subdivision of \([a, b]\) into partial intervals of sufficiently small maximal length so that the sum of the lengths of the intervals of this subdivision, where the oscillation of f is \(>\alpha\) is arbitrarily small. He shows as well that this condition is satisfied, not only for continuous and piecewise monotonic functions, but also for functions that may have an everywhere dense set of points of discontinuity. \(^{5}\)

The memoir of Riemann was not published until after his death, in 1867. But this time, the times were more favourable for this kind of research, and the ``Riemann integral'' took its place naturally in the stream of ideas that was leading then to a sustained study of the ``continuum'' and of functions of real variables (Weierstrass, Du Bois-Reymond, Hankel, Dini) and was to end up, with Cantor, at the blossoming of set theory. The form given by Riemann to the integrability condition suggested the idea of ``measure'' for the set of points of discontinuity of a function in an interval; but nearly 30 years were to pass by before a fruitful and convenient definition of this notion would be given.

The first attempts in this direction are due to Stolz, Harnack and Cantor (1884-1885); the first two, in order to define the ``measure'' of a bounded subset E of R, consider sets \(F \supset E\) which are finite unions of intervals, take for each F the sum of the lengths of the corresponding intervals, and say that the ``measure'' of E is the lower bound of these numbers; whereas Cantor, starting from the beginning in the space \(R^{n}\) , considers for a bounded set E and for \(\rho > 0\) the neighbourhood \(V(\rho)\) of E constructed of points whose distance to E is \(\leq \rho\) , and takes the lower bound of the ``volume'' of \(V(\rho)\) .⁶ With this definition, the ``measure'' of a set is equal to that of its closure, whence in particular it is possible that the ``measure'' of the union of two sets with no point in common can be strictly less that the sum of the ``measures'' of these two sets. No doubt in order to alleviate this last difficulty, Peano {[}246 a{]} and Jordan ({[}174 c{]}, v. I, pp.~28-31), a few years later, introduce, alongside the ``measure'' of Cantor \(\mu(A)\) of a set A contained in a block I, its ``interior measure'' \(\mu(I) - \mu(I - A)\) , and call ``measurable'' the sets A (now called ``quarriable'') for which these two numbers coincide. The union of two quarriable sets A, B with no point in common is then quarriable and has as ``measure'' the sum of the ``measures'' of A and B; but an open bounded set is not necessarily quarriable, and neither is the set of rational numbers contained in a bounded interval, which removed much of the interest in the notion of Peano-Jordan.

It is to E. Borel {[}32 a{]} that the credit for having known how to discern the faults of the previous definitions and for having seen how it was possible to remedy them is due. It was known since Cantor that every open set U in R is the union of a countable family of its ``components'', pairwise disjoint open intervals; instead of seeking to approach U ``from the outside'' by enclosing it in a finite sequence of intervals, Borel, relying on the preceding result, proposes taking as the measure of U (when U is bounded) the sum of the lengths of its components. Then he describes very summarily \(^{7}\) the class of sets (also called ``Borel'') that can be obtained, starting from the open sets, by iterating indefinitely the operations of countable union and of ``difference'' A - B, and indicates that, for these sets, a measure can be defined that possesses the fundamental property of complete additivity: if a sequence \((A_{n})\) is made up of pairwise disjoint Borel sets, the measure of their union (assumed bounded) is equal to the sum of their measures.

This definition was to inaugurate a new era in Analysis: on the one hand, in tandem with the contemporary work of Baire, it became the point of departure for a whole series of researches of a topological nature on the classification of sets of points (see p.~165); and above all, it was to be of use as the basis of the extension of the notion of integral, created by Lebesgue in the first years of the XXth century.

In his thesis {[}196 a{]}, Lebesgue starts by making precise and developing the succinct indications of E. Borel; imitating the method of Peano-Jordan, the ``exterior measure'' of a bounded set \(A \subset \mathbf{R}\) is defined as the lower bound of the measures of the open sets containing \(A\) ; then, if \(I\) is an interval containing \(A\) , the ``interior measure'' of \(A\) is the difference between the exterior measures of \(I\) and of \(I - A\) ; thus is obtained a notion of a ``measurable set'' which only differs from the initial ``constructive'' definition of Borel by the adjunction of a subset of a set of measure zero in the sense of Borel. This definition extended immediately to the spaces \(\mathbf{R}^n\) ; the old conception of the definite integral \(\int_{a}^{b} f(t) dt\) of a bounded function \(\geq 0\) as the ``area'' bounded by the curve \(y = f(x)\) , the straight lines \(x = a, x = b\) and \(y = 0\) , supplied thus an immediate extension of the Riemann integral to all the functions \(f\) for which the measure of the preceding set was defined. But the originality of

Lebesgue does not reside so much in the idea of this extension, \(^{8}\) than in his discovery of the fundamental theorem on the passage to the limit in the integral thus conceived, a theorem which appears in his work as a consequence of the complete additivity of the measure; \(^{9}\) he realises immediately all its importance, and makes it the key stone of the didactic exposé of his theory that he gives, already in 1904, in the famous ``Lessons on integration and the quest for primitive functions'' {[}196 c{]}. \(^{10}\)

We cannot describe here in detail the innumerable advances that the results of Lebesgue were to bring about in the study of the classical problems of infinitesimal Calculus. Lebesgue himself had already, in his thesis, applied his theory to the extension of the classical notions of length and area to more general sets than the usual curves and surfaces; concerning the considerable development of this theory over half a century, we send the reader back to the recent exposé by L. Cesari {[}59{]}. Let us mention also the applications to trigonometric series, developed by Lebesgue almost immediately after his thesis {[}196 b{]}, and which were to open new horizons for this theory, whose exploration is far from being over in our day (see {[}345{]}). Finally and especially, the definition of the spaces \(L^p\) and the Fischer-Riesz theorem ({[}111{]}, {[}260 a and c{]}; cf.~p.~213) brought to light the role that could be played in Functional Analysis by the new notion of integral; a role that could only grow with the subsequent generalisations of this notion, of which we shall speak in a moment.

Beforehand, we will stop a little longer with one of the problems to which Lebesgue devoted the most effort, the link between the notions of integral and of primitive. With the generalisation of the integral introduced by Riemann the question that arose naturally was to know if the classical correspondence between integral and primitive, valid for continuous functions, still subsisted in the more general case. But it is easy to give examples of functions f, integrable in the Riemann sense, and such that \(\int_{a}^{x}f(t)dt\) does not have a derivative (nor even a derivative on the right or on the left) at certain points;

conversely, Volterra had shown, in 1881, that a function \(F(x)\) may have a bounded derivative in an interval \(I\) , but be non integrable (in the Riemann sense) in \(I\) . By an analysis of great subtlety (where the theorem on the passage to the limit in the integral is far from being sufficient), Lebesgue succeeds in showing that, if \(f\) is integrable (in his sense) in \([a,b]\) , \(F(x) = \int_{a}^{x} f(t) dt\) has almost everywhere a derivative equal to \(f(x)\) {[}196 c{]}. Conversely, if a function \(g\) is differentiable in \([a,b]\) and if its derivative \(g' = f\) is bounded, \(f\) is integrable and we have the formula \(g(x) - g(a) = \int_{a}^{x} f(t) dt\) . But Lebesgue notes that the problem is much more complex when \(g'\) is not bounded; \(g'\) is not necessarily integrable in this case, and the first problem is therefore to characterise the continuous functions \(g\) for which \(g'\) exists almost everywhere and is integrable. By limiting himself to the case where one of the derived numbers \(^{11}\) of \(g\) is everywhere finite, Lebesgue showed that \(g\) is necessarily a function with bounded variation. \(^{12}\) Finally he establishes the converse of this last result: a function with bounded variation \(g\) has almost everywhere a derivative, and \(g'\) is integrable; but it is no longer necessarily true that

\[
g (x) - g (a) = \int_ {a} ^ {x} g ^ {\prime} (t) d t;\tag{22.2}
\]

the difference between the two parts of this relation is a function of bounded variation which is not constant and with zero derivative almost everywhere (a ``singular'' function). It remained to characterise the functions with bounded variation g such that the relation (22.2) held. Lebesgue established that these functions (called ``absolutely continuous'' by Vitali, who made a detailed study of them) are those that have the following property: the total variation of g in an open set U (sum of the total variations of g in each of the connected components of U) tends to 0 with the measure of U.

We will see below how, in a weakened form, these results would later acquire a much more general scope. In their initial form, their field of application remained fairly restricted, and did not go beyond the framework of the ``fine'' theory of functions of a real variable, of which we will not study here the subsequent development; we will be content with mentioning the deep work of Denjoy and his followers and successors (Perron, de la Vallée-Poussin, Khintchine, Lusin, Banach, etc.); the reader will find a detailed exposé in the book of S. Saks {[}268{]}.

One of the essential advances brought about by the Lebesgue theory concerns multiple integrals. This notion was introduced about the middle of the XVIIIth century, and first in the form of an ``indefinite integral'' (by analogy with the theory of the integral of functions of a single variable, \(\int\int f(x,y)dx dy\) denotes a solution of the equation \(\frac{\partial^{2}z}{\partial x\partial y}=f(x,y)\) ); but already from 1770, Euler had a very clear conception of the double integral extended to a bounded domain (bounded by arcs of analytic curves), and wrote down correctly the formula for evaluating such an integral by means of two simple integrals in succession ({[}108 a{]}, (1), v. XVII, pp.~289-315). It was not difficult to justify this formula starting from ``Riemann sums'', so long as the function being integrated was continuous, and the domain of integration not too complicated; but as soon as one wanted to tackle more general cases, the procedure of Riemann ran into serious difficulties \((f(x,y)\) may be integrable in the Riemann sense, without \(\int dx\int f(x,y)dy\) having a meaning when the simple integrals are taken in the Riemann sense). These difficulties vanish when we pass to the definition of Lebesgue; already this latter had showed in his thesis that, when \(f(x,y)\) is a bounded ``Baire function'', it is the same for the functions \(y\mapsto f(x,y)\) (for all x) and \(x\mapsto\int f(x,y)dy\) , and we have the formula

\[
\int \int f (x, y) d x d y = \int d x \int f (x, y) d y\tag{22.3}
\]

(integral taken in the rectangle). A little later, Fubini {[}121{]} brought an important complement to this result by proving that, if it is only assumed that \(f\) is integrable, then the set of \(x\) such that \(y \mapsto f(x, y)\) is not integrable is of zero measure, which allowed the extension of the formula (22.3) to this case.

Finally, in 1910 {[}157 e{]}, Lebesgue tackles the extension to multiple integrals of his results on the derivatives of simple integrals. He is thus brought to associate with a function \(f\) , integrable on all compact subsets of \(\mathbf{R}^n\) , the function of a set \(F(E) = \int_E f(\mathbf{x}) dx\) , defined for all integrable subsets \(E\) of \(\mathbf{R}^n\) , which generalises the concept of ``indefinite integral''; and he observes on this occasion that this function possesses the following two properties: \(1^\circ\) it is completely additive; \(2^\circ\) it is ``absolutely continuous'' in the sense that \(F(E)\) tends to 0 with the measure of \(E\) . The essential part of the memoir of Lebesgue consists in proving the converse of this proposition. \(^{13}\) But he does not stop there, and in the same direction, signals the possibility of generalising the notion of function with bounded variation, by considering functions of a measurable set \(F(E)\) , completely additive and such that \(\sum_n |F(E_n)|\) remains bounded for every countable partition of \(E\) into measurable subsets \(E_n\) . And, if he limits himself in fact to not considering such functions except on the set of blocks of \(\mathbf{R}^n\) , it is quite clear that there only remained one step to take to end up with the general notion of measure that J. Radon is going to define in 1913, gathering into the same synthesis the Lebesgue integral and the Stieltjes integral, of which we must now speak.

In 1894, T. Stieltjes published, under the title ``Research on continuous fractions'' {[}297{]}, a very original memoir where, starting from a question which is very restricted in appearance, were to be found stated and solved, with a rare elegance, problems of a quite new nature in the theory of analytic functions and that of functions of a real variable. \(^{14}\) In order to represent the limit of a certain sequence of analytic functions, Stieltjes had been brought, among others, to introduce on the straight line, the concept of a positive ``mass distribution'', a familiar notion for a long time in the physical sciences, but which until then had never been considered in mathematics except under restrictive hypotheses (in general, the existence of a ``density'' at every point, varying in a continuous way); he remarks that giving such a distribution is equivalent to that of an increasing function \(\phi(x)\) which gives the total mass contained in the interval with ends 0 and x for x \textgreater{} 0, and this mass having changed sign for x \textless{} 0, the discontinuities of \(\phi\) corresponding to masses ``concentrated at a point''. \(^{15}\) Stieltjes then forms, for such a mass distribution in an interval \([a, b]\) , the ``Riemann sums'' \(\sum_{i} f(\xi_{i})(\phi(x_{i+1}) - \phi(x_{i}))\) and shows that, when f is continuous in \([a, b]\) , these sums tend to a limit that he denotes \(\int_{a}^{b} f(x) d\phi(x)\) . Only needing to integrate continuous functions (and even differentiable functions) Stieltjes did not push further forward with the study of this integral \(^{16}\) and for about ten years this notion does not seem to have attracted attention. \(^{17}\) But in 1909, F. Riesz {[}260 d{]}, solving a problem stated a few years earlier by Hadamard (cf.~p.~213), proved that the Stieltjes integrals \(f \mapsto \int_{a}^{b} fd\phi\) are the most general continuous linear functionals on the space \(\mathcal{C}(I)\) of numerical functions continuous on \(I = [a, b] (\mathcal{C}(I))\) being furnished with the uniform convergence topology \(^{18}\) ; and the elegance and simplicity of this result caused almost immediately several generalisations of it. The most successful was that of J. Radon, in 1913 {[}256{]}: combining the ideas of F. Riesz and of Lebesgue, he showed how an integral could be defined by the procedures of Lebesgue, starting with an arbitrary ``completely additive function of a set'' (defined on measurable sets for the Lebesgue measure) instead of starting with the Lebesgue measure. In the notion of ``Radon measure'' on \(R^{n}\) , thus defined, that of function with ``bounded variation'' is found to be absorbed: the decomposition of such a function as a difference of two increasing functions is a particular case of the decomposition of a measure as a difference of two positive measures; in the same way, the ``measure with base \(\mu\) '' corresponds to the notion of an ``absolutely continuous'' function, and the decomposition of an arbitrary measure into a measure with base \(\mu\) and a measure foreign to \(\mu\) , to the Lebesgue decomposition of a function with bounded variation as the sum of an absolutely continuous function and a ``singular'' function. As well, Radon showed that the ``density'' with respect to \(\mu\) of a measure with base \(\mu\) exists still when \(\mu\) is a measure having as base the Lebesgue measure, by using a previous idea of F. Riesz (taken up again and popularised later by J. von Neumann among others), which consists in constructing an image of the measure \(\mu\) by a map \(\theta\) of \(R^{n}\) into R, chosen so that \(\theta(\mu)\) is the Lebesgue measure on R.

Almost immediately after the appearance of the memoir of Radon, Fréchet remarked that nearly all the results of this work could be extended to the case where the ``completely additive functions of a set'', instead of being defined on the measurable sets of \(R^{n}\) , are defined for certain subsets of an arbitrary set E (these subsets being such that the operations of countable union and of ``difference'' still give sets for which the function is defined). However, the expression of a measure with base \(\mu\) in the form g. \(\mu\) relied, with Lebesgue and Radon, on arguments using in an essential way the topology of \(R^{n}\) (and we have seen that the proof of Radon only applies if \(\mu\) is a measure with base the Lebesgue measure); it is only in 1930 that O. Nikodym {[}235{]} obtained this theorem in its general form, by a direct argument (notably simplified a few years later by J. von Neumann, thanks to the use of the properties of the spaces \(L^{2}\) ({[}324 g{]}, pp.~127-130)).

With the memoir of Radon, the general theory of integration could be considered as completed in its main lines; as subsequent substantial acquisitions, it is only possible to mention the definition of the infinite product of measures, due to Daniell {[}76 b{]}, and that of the integral of a function with values in a Banach space, given by Bochner in 1933 {[}24 a{]}, and which was a prelude to the more general study of the ``weak integral'' developed a few years later by Gelfand {[}125 a{]}, Dunford and Pettis ({[}95{]} and {[}96{]}). But it remained to popularise the new theory, and to make of it a mathematical instrument in general use, while the majority of mathematicians, around 1910, still saw in the ``Lebesgue integral'' only an instrument of great precision, delicate handling, destined only for researches of extreme subtlety and of an extreme abstraction. That was the work of Carathéodory, in a book which remained a classic for a long time {[}49{]} and which enriched Radon's theory with numerous original remarks.

But it is with this book also that the notion of integral, which had been in the front line of the concerns of Lebesgue (as was sufficiently shown by the titles of his thesis {[}196 a{]} and of his main work on these questions {[}196 c{]}) yields the field for the first time to that of measure, which had been for Lebesgue (as before him with Jordan) an auxiliary technical means. This change in point of view was due no doubt, with Carathéodory, to the excessive importance that he seems to have attached to ``p-dimensional measures''. \(^{19}\) Since then, the authors that have treated integration are divided between these two points of view, not without involving themselves in debates which have caused a lot of ink to flow if not a lot of blood. One group have followed Carathéodory; in their exposés ever more abstract and more axiomatised, measure, with all the technical refinements to which it lends itself, not only plays the dominant role, but also tends to lose contact with the topological structures to which it is in fact linked in the majority of problems where it is involved. Other exposés follow more or less closely a method already indicated in 1911 by W. H. Young, in a memoir unfortunately little noticed {[}339 b{]}, and developed since by Daniell. The first, dealing with the Lebesgue integral, parted company with the ``Cauchy integral'' of continuous functions with compact support, assumed known, in order to define successively the upper integral of lower semi-continuous functions, then of arbitrary numerical functions, whence a definition of integrable functions, modelled on that of Lebesgue for sets, by purely ``functional'' means. Daniell, in 1918 ({[}76 a{]}; cf.~{[}207{]}) extended this exposé, with some variants, to functions defined on an arbitrary set; in the same stream of ideas (and in close link with the methods used in spectral theory before Gelfand), we must also point out the memoir of F. Riesz {[}260 h{]} which puts into an elegant and concise form the few results of the theory of ordered spaces that play a role in the theory of Integration.

Rather than in works of exposition, more or less pleasant to read, but whose substantial content could not vary much, it is on the side of applications that progress arising from the theory of Integration since 1920 must be sought: the theory of probability (at other times the pretext for guessing and paradoxes, and become a branch of the theory of Integration since its

axiomatisation by Kolmogoroff {[}183{]}, but an autonomous branch with its own methods and problems); ergodic theory; spectral theory and harmonic analysis, since the discovery by Haar of the measure that bears his name, and the movement of ideas provoked by this discovery, have made of the integral one of the most important tools in the theory of groups.
